\documentclass[11pt]{article}
\usepackage[margin=1in]{geometry}
\usepackage{amsmath}
\usepackage{booktabs}
\usepackage{graphicx}
\usepackage{natbib}
\usepackage{cleveref}

\graphicspath{{figures/}}

\title{Mixed Academic Paper Stress Case}
\author{Test Corpus}
\date{}

\begin{document}
\maketitle

\begin{abstract}
This paper-style fixture combines numbered sections, cross-references, equations,
footnotes, a raster figure, a booktabs table, and a BibTeX bibliography. The
content is synthetic, but its structure resembles a compact empirical manuscript.
\end{abstract}

\section{Introduction}
\label{sec:introduction}

Academic conversion workflows are usually evaluated on a narrow set of examples.
This fixture uses a broader paper form in which prose, references, and floating
elements are interleaved across several pages. The design follows common article
conventions because many Word submissions retain the logical structure of the
LaTeX source while reflowing the exact page breaks.

The motivating setting is a multi-site program that introduced a scheduling
protocol across offices. Sites began with different administrative routines, which
creates variation in baseline completion and follow-up. The synthetic data in
this document have no empirical claim, although they make the source realistic
enough to exercise citation, footnote, and cross-reference handling. The notation
in \cref{sec:model} is intentionally simple so that equation conversion can be
checked without a specialized mathematical context.

Prior work on empirical design emphasizes clear descriptions of assignment and
measurement \citep{angrist2009mostly}. Reporting guidance also stresses that
tables and figures should remain interpretable when documents are moved between
formats \citep{gentzkow2014code}. This point matters for collaborative editing
because co-authors often review the Word file rather than the compiled PDF.

The introduction includes a real footnote to verify that body footnotes survive
the conversion path.\footnote{This is a body footnote with ordinary prose and an
inline expression, $n=48$.} A second footnote appears later in the paper so that
numbering can be inspected across sections.

\section{Research Design}
\label{sec:design}

The program introduced a standardized appointment sequence after an initial
baseline period. Offices retained discretion over local staffing, but the
sequence specified when reminders should be sent and when supervisors should
review delayed cases. The resulting design is presented as a compact panel of
site-period observations. The fixture uses this language because it creates
natural opportunities for equations, tables, and figure references in the same
document.

\begin{table}[htbp]
\centering
\caption{Baseline balance across synthetic sites}
\label{tab:balance}
\begin{tabular}{lccc}
\toprule
Measure & Early sites & Later sites & Difference \\
\midrule
Average weekly cases & 42.8 & 41.3 & 1.5 \\
Completion rate & 0.64 & 0.62 & 0.02 \\
Supervisor review rate & 0.31 & 0.34 & -0.03 \\
Missing contact share & 0.08 & 0.09 & -0.01 \\
\bottomrule
\end{tabular}
\end{table}

\cref{tab:balance} reports baseline characteristics for the synthetic site
groups. The table is small, but it includes the rule structure that often matters
for journal-formatted empirical papers. The columns contain both counts and
proportions so that decimal formatting appears alongside text labels.

The figure in \cref{fig:workflow} summarizes the assumed process. It is stored as
a local PNG to test figure resolution, media embedding, and caption placement.

\begin{figure}[htbp]
\centering
\includegraphics[width=0.74\linewidth]{workflow.png}
\caption{Synthetic scheduling workflow}
\label{fig:workflow}
\end{figure}

The process begins with an intake record, moves to a reminder stage, and ends
with either completion or supervisory review. The figure provides a deliberately
simple visual target so that discrepancies in scaling or placement are easy to
notice in side-by-side PDF inspection.

\section{Model}
\label{sec:model}

Let $Y_{it}$ denote the completion rate for site $i$ in period $t$. The main
specification is
\begin{equation}
Y_{it}=\alpha_i+\gamma_t+\beta D_{it}+X_{it}^{\top}\delta+\varepsilon_{it},
\label{eq:objective}
\end{equation}
where $D_{it}$ indicates exposure to the standardized sequence. The coefficient
$\beta$ is the parameter of interest in \cref{eq:objective}. The specification is
included to cover editable display equations and cross-references produced by
cleveref.

For descriptive reporting, the fixture also defines a normalized index:
\begin{align}
Z_{it}
  &= \frac{Y_{it}-\bar{Y}_{0}}{s_{0}},\\
\widehat{\Delta}
  &= \frac{1}{N_1}\sum_{i,t:D_{it}=1} Z_{it}
     - \frac{1}{N_0}\sum_{i,t:D_{it}=0} Z_{it}.
\end{align}
These equations create multi-line mathematical content while remaining close to
ordinary applied notation.

\section{Results}
\label{sec:results}

The estimated values are intentionally moderate. A short discussion follows to
force the document beyond a single page and to test page-level rendering in the
harness. The fixture refers back to \cref{sec:design}, \cref{tab:balance}, and
\cref{fig:workflow} so that the auxiliary file contains references of several
types.

The first result is that completion increases after the standardized sequence is
introduced. The second result is that supervisory review becomes more frequent in
sites with larger initial backlogs, which is consistent with the intended
workflow. These statements are synthetic and should be read only as conversion
content. Their value for the test suite lies in the density of ordinary academic
structure around them.\footnote{The second body footnote appears after the table
and figure references, allowing later footnote numbering to be checked.}

Additional prose extends the fixture to several pages. The conversion path should
preserve section headings, ordinary paragraphs, citations, footnotes, equation
numbers, and captions even when the Word output reflows differently from the PDF.
The compiled file remains the visual reference, while the Word document provides
an editable representation suitable for review.

The paper form also exercises bibliography placement. The final section cites a
methods reference again \citep{imbens2015causal} and then prints the bibliography
through standard BibTeX commands. The combination of \verb|\cref| commands and
natbib citations gives the converter a mixed reference workload.

\section{Conclusion}
\label{sec:conclusion}

This mixed fixture is meant to be ordinary rather than exotic. It contains enough
structure to expose failures in floats, references, footnotes, equations, and
bibliography handling, while each individual element remains small enough for
manual inspection. The harness records page counts and conversion summaries so
that repeated runs can be compared without changing the source files.

The conclusion adds another paragraph to stabilize page length. Reviewers can
compare \cref{eq:objective}, \cref{tab:balance}, and \cref{fig:workflow} between
the compiled PDF and the rendered Word PDF. Differences in exact pagination are
expected in a reflowed Word document, but missing elements or unresolved commands
should be visible in the rasterized outputs.

\bibliographystyle{plainnat}
\bibliography{refs}

\end{document}
